As an additional approach we explore a greedy algorithm with two goals in mind:
to investigate whether this simpler approach performs similarly,
and to verify that optimal rulesets for simple cases (low $C$ and $D$)
are similar to those produced by our earlier, more complicated approach.

In the greedy approach, we first find through brute force complete enumeration a rule with minimum \zo loss,
so with maximum train accuracy.
Then, the positively classified points are removed from the dataset and the cycle repeats.
Note that this greedy method is essentially a strategy to find a rule list rather than a rule set.
For $C=1$, the greedy approach finds an optimal ruleset.

Mathematically, we consider in step $\tau$ the points $\I_\tau$ and the corresponding positive points $\P_\tau$ and negative points $\N_\tau.$
Denote by rule $k_\tau$ the rule found in step $\tau$.
Then $\I_1 = \I$ and
\[\I_{\tau + 1} = \I_\tau \setminus \{i \in \I_\tau \fw k_\tau \in \K_{i}\} = \I_\tau \setminus \I_{\{k_\tau\}}, \tquad \tau \in \mathbb{N}_+.\]
Equivalently,
\[\I_{\tau + 1} = \{i \in \I_\tau \fw y_{\{k_\tau\}}(\bm{X}_i) \neq 1\}, \tquad \tau \in \mathbb{N}_+.\]

\begin{algorithm}[h!]
\caption{Greedy Ruleset Selection}
\label{alg:greedy}
\begin{algorithmic}[1]
\Require $\I$, $\bm{X}$, $\bm{y}$, $C$, $D$
\State $\I_1 \gets \I$
\For{$\tau = 1, \dots, C$}
    \State Find rule $k_\tau$ with at most $D$ clauses minimizing \zo loss w.r.t. $\I_\tau,$ $\bm{X},$ $y$
    \State $\I_{\tau+1} \gets \{i \in \I_{\tau} \fw \hat{y}_{\{k_\tau\}}(\bm{X}_i) \neq 1\}$
\EndFor \\
\Return $(k_\tau)_{\tau=1,\dots,C}$
\end{algorithmic}
\end{algorithm}

An advantage of the greedy algorithm is that it is easy to understand and to implement.
Furthermore, discretization via thresholding is done in every step separately.
This means that new, more relevant threshold values may be found in later steps,
when a part of the points is no longer considered.
A disadvantage is that there is no optimality guarantee for $C>1$.
Additionally, the Python implementation, while highly readable, is slow.
This is in part due to enumeration of all rules in each step, which is slow for larger values of $D$
(in our experiments, for $D\geq 3$).
Another aspect is the high-level nature of Python, hurting performance, despite the use of NumPy where possible.

The results for $D=1$ are summarized in \autoref{tab:greedy_yq_4_5_mw_0_C_3_D_1}.
While the greedy method was configured to continue after step two, there did not exist a rule that improved
performance.
The first rule provides reasonable performance.
The second rule provides barely any improvement.
For both $Q_y=0.5$ and $Q_y=0.6$, the same features appear in the rules.

The results for $D=2$ are summarized in \autoref{tab:greedy_yq_4_5_mw_0_C_3_D_2}.
In both tables, the value of adding another rule quickly diminishes.
The second and third rules barely increase accuracy.
The third rule found covers only a small fraction of the remaining positive points.
%Compared to \autoref{tab:threshold_interpretable_results_C_1} and \autoref{tab:threshold_interpretable_results_C_2},
%describing results of our earlier approach,
%the accuracy observed is similar.
Qualitatively, there is some overlap in the rules found, but there are differences.

The previous method (\autoref{sec:threshold-results}) performs similar to the greedy method.
Since for $C=1$ greedy is optimal, this suggests that the previous method is working as intended.
Sometimes, the greedy method outperforms the previous method.
There are two contributing factors for this.
First, the greedy method has more freedom in later steps.
Second, the previous approach is a heuristic, as it uses price-then-branch, and it was configured to use
only the heuristic pricer.
Nevertheless, overall performance is similar.

The conclusion of the experiments in this chapter is two-fold.
\begin{itemize}
    \item Compared to \autoref{sec:boolean-results}, decreasing complexity only moderately decreases performance.
    Instead, very simple rulesets can achieve a comparable accuracy, even when only using one clause per rule,
    which was not expected.
    It appears that the data is highly correlated and noisy, making it unsuitable for this off-the-shelf classifier model.
    \item The rulesets found are not consistent.
    The cross-validation experiment showed that rulesets vary among splits.
    The inconsistency of rules is supported by the wide variety of TF combinations observed in this chapter's tables.
\end{itemize}
We will therefore alter the model in an attempt to make it more robust and handle noise better, hopefully
providing more meaningful rulesets.

%\begin{table}
%    \centering
%    \bwtable{gen/PDL1ex2/tab_greedy_yq_4_mw_0_C_3_D_2}
%    \caption{}
%    \label{tab:greedy_yq_4_mw_0_C_3_D_2}
%\end{table}
%
%\begin{table}
%    \centering
%    \bwtable{gen/PDL1ex2/tab_greedy_yq_5_mw_0_C_3_D_2}
%    \caption{}
%    \label{tab:greedy_yq_5_mw_0_C_3_D_2}
%\end{table}

\begin{table}
    \centering
    \bwtable{gen/PDL1ex2/tab_greedy_yq_4_mw_0_C_3_D_1}
    \vspace{1em}
    \bwtable{gen/PDL1ex2/tab_greedy_yq_5_mw_0_C_3_D_1}
    \caption{Rules found in each step of the greedy algorithm, for $Q_y=0.5$ (top) and $Q_y=0.6$ (bottom).
    The full Broad dataset is used for training and $D=1$.
    In both cases, accuracy cannot be improved by any additional rule after step two.}
    \label{tab:greedy_yq_4_5_mw_0_C_3_D_1}
\end{table}

\begin{table}
    \centering
    \bwtable{gen/PDL1ex2/tab_greedy_yq_4_mw_0_C_3_D_2}
    \vspace{1em}
    \bwtable{gen/PDL1ex2/tab_greedy_yq_5_mw_0_C_3_D_2}
    \caption{Rules found in each step of the greedy algorithm, for $Q_y=0.5$ (top) and $Q_y=0.6$ (bottom).
    The full Broad dataset is used for training and $D=2$.}
    \label{tab:greedy_yq_4_5_mw_0_C_3_D_2}
\end{table}


%It can be formulated as follows:~
%\begin{subequations}
%    \label{eq:TP}
%    \begin{align}
%        \min \quad & \sum_{i \in \cN} \delta_i - \sum_{i \in \cP} \mu_i \delta_i +
%        \lambda \label{eq:TP-obj} \\
%        \st \quad & |J|\delta_i + \sum_{j \in \J} z_{ij} \leq |J|, & i \in \I, \label{eq:TP-con1}\\
%        & \delta_i +  \sum_{j \in J} z_{ij} \geq  1, & i \in \I,  \label{eq:TP-con2}\\
%        & z_{ij} \geq z_{i'j} & \forall\, i \in \I,\, j \in \J, \, :\, X_{ij} \leq X_{i'j}, \label{eq:TP-con3} \\
%        & \mathrlap{\delta\in \{0,1\}^\I,~z\in\{0,1\}^{\I\times \J}.}\label{eq:TP-vars}
%    \end{align}
%\end{subequations}
%Here $z_{ij}$ encodes whether point $i$ is cut off by feature $j$ in the chosen rule.
%Note that the rule complexity is not capped at $D$ in this formulation.
%In practice, this IP is hard to solve, due to its size.

